\documentclass[11pt,a4paper]{article}

\usepackage[utf8]{inputenc}
\usepackage[T1]{fontenc}
\usepackage[margin=2.5cm]{geometry}
\usepackage{amsmath}
\usepackage{booktabs}
\usepackage{xcolor}
\usepackage{listings}
\usepackage{hyperref}

\lstset{
  language=C,
  basicstyle=\ttfamily\small,
  keywordstyle=\color{blue!70!black}\bfseries,
  commentstyle=\color{green!40!black}\itshape,
  stringstyle=\color{red!60!black},
  numbers=left,
  numberstyle=\tiny\color{gray},
  frame=single,
  breaklines=true,
  columns=fullflexible,
  morekeywords={int32_t,int64_t,uint16_t,fxpt_4_28,rgb565}
}

\title{CS-473 System Programming for Systems on Chip\\[0.3em]
       \large Practical Work 2 -- The Mandelbrot Set\\
       Part 1: Fixed-Point Implementation}
% TODO: add your partner's name
\author{Matheo Chacon}
\date{\today}

\begin{document}
\maketitle

\section{Introduction}

The reference program \texttt{fractal\_flpt} computes the Mandelbrot set with
the C \texttt{float} type. The OpenRISC processor of the virtual prototype has
no floating-point unit, so every \texttt{float} addition and multiplication is
done by a software library, which is slow. In this first part we rewrote the
program as \texttt{fractal\_fxpt}, in which every coordinate and intermediate
value is a 32-bit fixed-point integer. This report explains the format we
chose and the changes we made to the three source files.

\section{Choice of the fixed-point format}

The Mandelbrot iteration is
\[
  z_{n+1} = z_n^2 + c, \qquad z_0 = 0,
\]
which, written with real and imaginary parts, becomes
\[
  x_{n+1} = x_n^2 - y_n^2 + c_x, \qquad y_{n+1} = 2x_ny_n + c_y .
\]
A point escapes as soon as $x_n^2 + y_n^2 \geq 4$.

To choose the number of integer bits, we need the largest value a 32-bit
variable has to hold:
\begin{itemize}
  \item With the default view, $c_x \in [-2, 1]$ and $c_y \in [-1.5, 1.5]$.
  \item We only compute $z_{n+1}$ while $|z_n|^2 < 4$, so
        $|x_n^2 - y_n^2| < 4$ and $|2x_ny_n| \leq x_n^2 + y_n^2 < 4$.
  \item As a result, $|x_{n+1}| < 4 + 2 = 6$ and $|y_{n+1}| < 4 + 1.5 = 5.5$.
\end{itemize}
The integer part must therefore hold values up to about $\pm 6$, which takes
one sign bit and three magnitude bits. We put all remaining bits in the
fractional part and chose the \textbf{Q4.28} format:

\begin{center}
\begin{tabular}{@{}ll@{}}
  \toprule
  Range      & $[-8,\; 8 - 2^{-28}]$ \\
  Resolution & $2^{-28} \approx 3.7 \times 10^{-9}$ \\
  Constants  & $4.0 = 2^{30}$, \quad $3.0 = 3 \cdot 2^{28}$ \\
  \bottomrule
\end{tabular}
\end{center}

With only three integer bits (Q3.29) the range would be $[-4, 4)$, which is
too small for the intermediate values computed above. Choosing exactly as many
integer bits as needed leaves the largest possible number of fractional bits.
This matters for the next exercise: when we zoom in, the pixel step
\texttt{delta} becomes very small, and 28 fractional bits allow deep zooms
before neighbouring pixels map to the same coordinate.

The format is declared in \texttt{include/fractal\_fxpt.h}:
\begin{lstlisting}
#define DECIMAL_BITS 28
#define FIXED_MAX ((1 << (32 - DECIMAL_BITS)) - 1)
#define FIXED_MIN (-(1 << (32 - DECIMAL_BITS)))
#define DOUBLE_TO_FXPT_4_28(a)                                     \
  ((fxpt_4_28)(((a) >= (FIXED_MAX))   ? (INT32_MAX)                \
               : ((a) <= (FIXED_MIN)) ? (INT32_MIN)                \
                                      : ((a) * (1 << DECIMAL_BITS))))

typedef int32_t fxpt_4_28; // Q4.28
\end{lstlisting}

The macro \texttt{DOUBLE\_TO\_FXPT\_4\_28} converts constants such as
\texttt{-2.0} into Q4.28 by multiplying them by $2^{28}$, and saturates any
value outside the range. Because it is only applied to constant expressions,
the compiler evaluates it at compile time, so the program never uses
floating-point arithmetic at run time.

\section{Implementation}

\subsection{Header: changing the function signatures}

Every \texttt{float} parameter in the header was replaced by
\texttt{fxpt\_4\_28}:
\begin{lstlisting}
typedef uint16_t (*calc_frac_point_p)(fxpt_4_28 cx, fxpt_4_28 cy,
                                      uint16_t n_max);

uint16_t calc_mandelbrot_point_soft(fxpt_4_28 cx, fxpt_4_28 cy,
                                    uint16_t n_max);

void draw_fractal(rgb565 *fbuf, int width, int height,
                  calc_frac_point_p cfp_p, iter_to_colour_p i2c_p,
                  fxpt_4_28 cx_0, fxpt_4_28 cy_0,
                  fxpt_4_28 delta, uint16_t n_max);
\end{lstlisting}
As the handout specifies, we left the colour-mapping functions
(\texttt{iter\_to\_bw}, \texttt{iter\_to\_grayscale},
\texttt{iter\_to\_colour}, \ldots) and \texttt{ilog2} unchanged, since they
only work on iteration counts.

\subsection{Multiplication (hint 1)}

When two Q4.28 numbers are multiplied, the exact result has 64 bits in the
Q8.56 format: the integer bits add up ($4+4$) and so do the fractional bits
($28+28$). To get back to Q4.28, we compute the product in 64 bits and shift
it right by 28:
\[
  a \cdot b \;\longrightarrow\; \bigl((\text{int64})\,a \cdot (\text{int64})\,b\bigr) \gg 28
  \quad (\text{Q8.56} \to \text{Q8.28}).
\]
Before the shift we cast both operands to \texttt{int64\_t}. Without this
cast, C would perform a 32-bit multiplication and discard the upper half of the
result. After the shift the value is in Q8.28 and stored in an
\texttt{int64\_t}, so even a square such as $x^2$ that exceeds 8 is
represented exactly.

To multiply by two in $2xy$, we shift left by one bit instead of performing a
multiplication.

\subsection{Overflow handling (hint 2)}

Integer arithmetic gives no warning when a result overflows: it silently wraps
around. We protect against overflow in three places:
\begin{enumerate}
  \item \textbf{Inside the iteration}, we keep every intermediate value
        ($x^2$, $y^2$, $xy$ and the new $x$ and $y$) in 64-bit variables.
        Before we store a new value back into a 32-bit \texttt{fxpt\_4\_28},
        we check that it fits. If it does not fit, the point is clearly
        diverging, so we stop iterating.
  \item \textbf{The escape test} $x^2 + y^2 \geq 4$ is also computed in
        64 bits, so the sum cannot overflow even though it may exceed 8.
  \item \textbf{When stepping through the pixels} in \texttt{draw\_fractal},
        we use \texttt{fixed\_add}, which saturates to
        \texttt{INT32\_MAX}/\texttt{INT32\_MIN} instead of wrapping around.
\end{enumerate}
In the default view, the bounds of Section~2 show that case~1 never
occurs, since all values stay below 6. The check is there for later zoom and
pan exercises, where $c$ can lie outside the default window.

\subsection{The Mandelbrot kernel}

The floating-point kernel was:
\begin{lstlisting}
float x = cx, y = cy;
do {
  xx = x * x;  yy = y * y;  two_xy = 2 * x * y;
  x = xx - yy + cx;
  y = two_xy + cy;
  ++n;
} while (((xx + yy) < 4) && (n < n_max));
\end{lstlisting}
Our fixed-point version is:
\begin{lstlisting}
uint16_t calc_mandelbrot_point_soft(fxpt_4_28 cx, fxpt_4_28 cy,
                                    uint16_t n_max) {
  fxpt_4_28 x = 0;
  fxpt_4_28 y = 0;
  uint16_t n = 0;
  const fxpt_4_28 limit = (fxpt_4_28)4 << DECIMAL_BITS;

  while (n < n_max) {
    int64_t xx_64 = ((int64_t)x * (int64_t)x) >> DECIMAL_BITS;
    int64_t yy_64 = ((int64_t)y * (int64_t)y) >> DECIMAL_BITS;

    if ((xx_64 + yy_64) >= limit) {
      break;
    }

    int64_t xy_64 = ((int64_t)x * (int64_t)y) >> DECIMAL_BITS;

    int64_t next_x = xx_64 - yy_64 + (int64_t)cx;
    int64_t next_y = (xy_64 << 1) + (int64_t)cy;

    if (next_x > INT32_MAX || next_x < INT32_MIN ||
        next_y > INT32_MAX || next_y < INT32_MIN) {
      break;
    }

    x = (fxpt_4_28)next_x;
    y = (fxpt_4_28)next_y;
    ++n;
  }
  return n;
}
\end{lstlisting}
The structure is slightly different from the original \texttt{do-while}
loop. We start from $z_0 = 0$ instead of $z_1 = c$, and we test for escape
\emph{before} computing the next value instead of after. Both loops stop at
the first $n$ for which $|z_n|^2 \geq 4$, so they return the same iteration
count. Our version also skips the cross product $xy$ for a point that has
already escaped.

\subsection{Drawing loop and \texttt{main}}

In \texttt{draw\_fractal}, the pixel coordinates are now
\texttt{fxpt\_4\_28} values, and they are advanced with the saturating adder:
\begin{lstlisting}
fxpt_4_28 fixed_add(fxpt_4_28 a, fxpt_4_28 b) {
  int64_t result = (int64_t)a + (int64_t)b;
  if (result > INT32_MAX) return INT32_MAX;
  if (result < INT32_MIN) return INT32_MIN;
  return (fxpt_4_28)result;
}
\end{lstlisting}

In \texttt{main\_fxpt.c}, the view-port constants are converted at compile
time:
\begin{lstlisting}
const fxpt_4_28 FRAC_WIDTH = DOUBLE_TO_FXPT_4_28(3.0);
const fxpt_4_28 CX_0       = DOUBLE_TO_FXPT_4_28(-2.0);
const fxpt_4_28 CY_0       = DOUBLE_TO_FXPT_4_28(-1.5);
...
fxpt_4_28 delta = FRAC_WIDTH / SCREEN_WIDTH;
\end{lstlisting}
Dividing a fixed-point number by a plain integer requires no scaling, so
\texttt{delta} is an ordinary integer division. In this case it is even
exact: $3 \cdot 2^{28} / 2^9 = 3 \cdot 2^{19}$.

We also made two smaller changes to \texttt{main}:
\begin{itemize}
  \item The frame buffer ($512 \times 512 \times 2$ bytes $= 512$\,KiB) was
        moved from the stack to a \texttt{static} global array. A buffer this
        large is too big for the stack.
  \item The screen size is now defined with \texttt{\#define} instead of
        \texttt{const int}. In C, a \texttt{const int} is not a constant
        expression, so it cannot be used to size a static array.
\end{itemize}

\section{Results and observations}

The program compiles without errors, and we ran it on the virtual prototype.
It draws the same Mandelbrot image as the floating-point version. With the
default view and $N_{max} = 64$, the 28 fractional bits are many more than
the pixel step needs ($\delta = 3/512 \approx 5.9 \times 10^{-3}$), so we
saw no visible difference between the two images.

% TODO: add the measured execution times / cycle counts if you have them, e.g.
% \begin{center}
% \begin{tabular}{lr}
%   \toprule
%   Version & Execution time \\
%   \midrule
%   \texttt{fractal\_flpt} (soft-float) & ... \\
%   \texttt{fractal\_fxpt} (Q4.28)      & ... \\
%   \bottomrule
% \end{tabular}
% \end{center}
The main difference is speed: the fixed-point version draws the image
noticeably faster than the floating-point version. The \texttt{float} version
calls a software library routine for every addition and multiplication. In the
fixed-point version, additions are single integer instructions, and each
multiplication is one 64-bit integer multiplication followed by a shift.

\section{Conclusion}

We converted the Mandelbrot program from \texttt{float} to the Q4.28
fixed-point format. Four integer bits are just enough to hold every
intermediate value of the iteration, and the remaining 28 bits give a fine
resolution for the zoom exercise that comes later. Computing products in
64 bits and shifting the result right by 28 keeps the decimal point in place,
and checking ranges before storing values back into 32 bits prevents silent
overflow. The result is an image identical to the floating-point version,
computed without any floating-point operation at run time.

\end{document}
